% 两枚公平骰子的三种记录方式。样本点(i,j)，i为第一枚、j为第二枚点数，36个等可能。
% (a) 完整记录：每个格子单独可辨，信息为全部2^36个子集；
% (b) 只记录点数和S=i+j：同一反对角线上的格子不可分辨，11个原子；S=7一条有6格，概率6/36；
% (c) 只记录和的奇偶：两个原子，各18格，概率1/2。同色块为不可再分的原子；(b)中红框标出S=7。
\begin{tikzpicture}[line width=.6pt,font=\fontsize{8}{10}\selectfont,>=Stealth]
\def\c{4mm}
% ---------- (a) ----------
\begin{scope}
 \foreach \i in {1,...,6}{\foreach \j in {1,...,6}{
   \draw[FigureStroke,fill=FigureBlueFill] ({(\i-1)*\c},{(\j-1)*\c}) rectangle ++(\c,\c);}}
 \foreach \k in {1,...,6}{
   \node[FigureMuted,font=\fontsize{7.5}{9}\selectfont] at ({(\k-.5)*\c},-2.2mm) {$\k$};
   \node[FigureMuted,font=\fontsize{7.5}{9}\selectfont] at (-2.2mm,{(\k-.5)*\c}) {$\k$};}
 \node[FigureInk] at (12mm,-6.5mm) {第一枚 $i$};
 \node[FigureInk,rotate=90] at (-6.5mm,12mm) {第二枚 $j$};
 \node[FigureInk] at (12mm,-12mm) {(a) 完整结果：36个原子};
\end{scope}
% ---------- (b) ----------
\begin{scope}[xshift=40mm]
 \foreach \i in {1,...,6}{\foreach \j in {1,...,6}{
   \pgfmathtruncatemacro{\s}{mod(\i+\j,3)}
   \ifcase\s \def\fc{FigureBlueFill}\or \def\fc{FigureGreenFill}\or \def\fc{FigureYellowFill}\fi
   \fill[\fc] ({(\i-1)*\c},{(\j-1)*\c}) rectangle ++(\c,\c);}}
 % light cell grid; colour bands (sum mod 3) mark the atoms {S=s}
 \draw[FigureStroke] (0,0) rectangle (24mm,24mm);
 \foreach \i in {1,...,5}{\foreach \j in {1,...,6}{
   \draw[FigureGrid,line width=.4pt] ({\i*\c},{(\j-1)*\c}) -- ++(0,\c);}}
 \foreach \i in {1,...,6}{\foreach \j in {1,...,5}{
   \draw[FigureGrid,line width=.4pt] ({(\i-1)*\c},{\j*\c}) -- ++(\c,0);}}
 \draw[FigureRed,line width=1.1pt,rounded corners=1pt] (0,24mm) -- (0,20mm) -- (4mm,20mm) -- (4mm,16mm) -- (8mm,16mm) -- (8mm,12mm) -- (12mm,12mm) -- (12mm,8mm) -- (16mm,8mm) -- (16mm,4mm) -- (20mm,4mm) -- (20mm,0) -- (24mm,0) -- (24mm,4mm) -- (20mm,4mm) -- (20mm,8mm) -- (16mm,8mm) -- (16mm,12mm) -- (12mm,12mm) -- (12mm,16mm) -- (8mm,16mm) -- (8mm,20mm) -- (4mm,20mm) -- (4mm,24mm) -- cycle;
 \node[FigureRed,anchor=west] at (25mm,2mm) {$S=7$};
 \node[FigureMuted,font=\fontsize{7.5}{9}\selectfont,anchor=west] at (25mm,22mm) {$S=12$};
 \node[FigureMuted,font=\fontsize{7.5}{9}\selectfont,anchor=east] at (-1mm,2mm) {$S=2$};
 \node[FigureInk] at (12mm,-6.5mm) {同一反对角线不可分};
 \node[FigureInk] at (12mm,-12mm) {(b) 点数和：11个原子};
\end{scope}
% ---------- (c) ----------
\begin{scope}[xshift=82mm]
 \foreach \i in {1,...,6}{\foreach \j in {1,...,6}{
   \pgfmathtruncatemacro{\s}{mod(\i+\j,2)}
   \ifcase\s \def\fc{FigureBlueFill}\or \def\fc{FigureYellowFill}\fi
   \fill[\fc] ({(\i-1)*\c},{(\j-1)*\c}) rectangle ++(\c,\c);
   \draw[FigureGrid,line width=.4pt] ({(\i-1)*\c},{(\j-1)*\c}) rectangle ++(\c,\c);}}
 \draw[FigureStroke] (0,0) rectangle (24mm,24mm);
 \node[draw=FigureStroke,fill=FigureBlueFill,minimum size=3mm,inner sep=0pt] (e) at (28mm,18mm) {};
 \node[FigureInk,anchor=west] at (29.8mm,18mm) {和为偶};
 \node[draw=FigureStroke,fill=FigureYellowFill,minimum size=3mm,inner sep=0pt] (o) at (28mm,12mm) {};
 \node[FigureInk,anchor=west] at (29.8mm,12mm) {和为奇};
 \node[FigureInk] at (12mm,-6.5mm) {同色格子不可分};
 \node[FigureInk] at (15mm,-12mm) {(c) 奇偶：2个原子};
\end{scope}
\end{tikzpicture}
